\documentclass[10pt, aspectratio=169]{beamer}
\usepackage[russian]{babel}
\usepackage{pgfpages}
\usepackage{graphicx}
\usepackage{tabularx}
\usepackage[T1]{fontenc}
\usepackage{tgbonum}

\setbeameroption{hide notes}
\setbeamertemplate{footline}{
    \leavevmode
    \begin{beamercolorbox}[ht=2.5ex,dp=1.125ex,leftskip=1em,rightskip=.3em]{author in head/foot}
        \usebeamerfont{author in head/footer}Алтунина Екатерина Юрьевна
        \hfill
        \usebeamerfont{page number in head/foot}\insertframenumber\,/\,\inserttotalframenumber\hspace{1em}
    \end{beamercolorbox}
    \vspace{0.2em}
}
% TITLE PAGE
\title[Short title]{Визуализация теней в реальном времени }
\date[02.02.2026]{Апрель, 2026}
\author[author]{Алтунина Екатерина Юрьевна}
\institute{Руководитель: Романов Алексей Алексеевич\and ИПКН ИТМО, программа «Разработка программного обеспечения»\and(Санкт-Петербург)}


% \setbeamerfont{normal text}{family=\sffamily}
% \setbeamerfont{structure}{family=\sffamily}

\begin{document}
%_____________________________________________________
% Title page
\frame[plain]{
\titlepage
}
%_____________________________________________________
% ToC
% \frame{{Содержание}
% \tableofcontents
% }
%_____________________________________________________
% 
\section{Section 1}
\subsection{Subsection 1.1}
\frame{{Введение в предметную область} 

\begin{columns}
    \column{0.65\textwidth}
    Тени - ключевой элемент любой системы визуализации: 
    они определяют расстояния между объектами, особенно в картографии.

    В условиях веб-браузера приходится находить разумный 
    компромисс между реалистичностью теней и скоростью работы сцены.
    \column{0.3\textwidth}
    \includegraphics[width=\textwidth]{no shadows.PNG}
    \includegraphics[width=\textwidth]{shadows.PNG}
\end{columns}
}

\subsection{Subsection 1.2}
\frame{{Существующие решения} 

\begin{table}[h]
\centering
\footnotesize
\begin{tabular}{|l|c|c|c|c|c|}
\hline
\textbf{Метод} & \textbf{Скорость} & \textbf{Память} & \textbf{Область} & \textbf{Качество} & \textbf{Вывод} \\
\hline
Планарные тени    & + & + & - & - & Малая обл. применения \\
\hline
Теневой объем     & {\tiny$\pm$} & {\tiny$\pm$} & + & + & Зависит от покрытых пикселей \\
\hline
Трассировка лучей & - & - & + & + & Высокая сложность \\
\hline
Карты теней       & {\tiny$\pm$} & {\tiny$\pm$} & + & {\tiny$\pm$} & \textbf{Используем} \\
\hline
\end{tabular}
\caption{Сравнение методов рендеринга теней}
\end{table}

}

\subsection{Subsection 1.3}
\frame{{Карты теней} 

Универсальное и высокопроизводительное решение.

\begin{itemize} 
\item (+) Подходит для создания теней от различных источников света (точечных, направленных, прожекторов)
\item (+) Относительно предсказуемы в плане оценки стоимости построения карты теней
\item (--) Наличие артефактов
\item (--) Требуется строить одну или несколько карт теней для каждого источника света
\end{itemize}

}

\subsection{Subsection 1.7}
\frame{{Аналитические тени} 

Аналитические тени -
это метод вычисления теней без использования карт глубины, текстур или растеризации. 
Вместо этого тень рассчитывается математически, по формуле, прямо в шейдере, на основе формы объекта-источника тени.
\vspace{0.5em}

\begin{center}
    \includegraphics[width=0.6\textwidth]{analytic-shadows.jpg}
\end{center}
}

%_____________________________________________________
%
\section{Section 2}
\subsection{Subsection 2.1}
\frame{{Цель и задачи} 
{} 

Цель: оценить применимость и производительность метода аналитических теней в браузере и скомбинировать его с классическими подходами теневых карт.
\vspace{1em}

Задачи:
\begin{itemize} 
    \item Создать Web-песочницу с использованием WebGPU, реализовать навигацию по сцене и классический подход к визуализации теней - карту теней с фиксированным шагом
    \item Разработать систему раздельного рендеринга динамических/статических объектов и аналитические тени для динамических объектов на простых геометрических формах. Предусмотрев возможность изменения параметров, влияющих на отрисовку теней.
    \item Совместить оба способа, оценить корректность отрисовки теней и провести анализ производительности (FPS, время кадра, нагрузка на GPU) при различных сценариях
    \item Провести валидацию полученных данных
\end{itemize}
}

\section{Section 3}
\subsection{Subsection 3.1}
\frame{{Создание веб-песочницы} 
{}

\begin{columns}
    \column{0.65\textwidth}
    На первом этапе разработки:
    \begin{itemize} 
        \item Реализована навигация по сцене
        \item Поддержано два вида камеры -  ортографическая и перспективная.
        \item Для работы с камерой и навигацией - Three.JS
    \end{itemize}
    \column{0.3\textwidth}
    \includegraphics[width=1.0\textwidth]{1.jpg}
\end{columns}

}

\section{Section 3}
\subsection{Subsection 3.2}
\frame{{Реализация метода каскадных теней} 
{}

На втором этапе разработки:
\begin{itemize} 
    \item Реализован метод каскадных теней (до 8 каскадов)
    \item Динамический расчет границ отсечения каскадов в зависимости от положения камеры
    \item Добавлена визуализация отладочного режима каскадов (отображение цветовых зон)
\end{itemize}

\begin{center}
    \includegraphics[width=0.6\textwidth]{casc.png}
\end{center}

}

\section{Section 3}
\subsection{Subsection 3.3}
\frame{{Борьба с артефактами} 
{}

\vspace{0.5em}
\begin{table}[h]
    \centering
    \begin{tabular}{ccc}
        \textbf{Алиасинг} & \textbf{«Теневые прыщи»} & \textbf{Обрезание теней} \\
        \includegraphics[width=0.2\textwidth]{aliasing-no-pcf.jpg} & 
        \includegraphics[width=0.3\textwidth]{shadow-acne.jpg} & 
        \includegraphics[width=0.3\textwidth]{cutoutshadows.jpg} \\
        \parbox{0.3\textwidth}{\centering \textbf{Метод борьбы:} PCF (Percentage Closer Filtering) --- Сглаживание краёв} & 
        \parbox{0.3\textwidth}{\centering \textbf{Метод борьбы:} Смещение (depth bias)} & 
        \parbox{0.3\textwidth}{\centering \textbf{Метод борьбы:} Вписывание сцены во фруструм} \\
    \end{tabular}
\end{table}
}


\section{Section 3}
\subsection{Subsection 3.5}
\frame{{Система для разделения рендеринга динамических и статических объектов} 
{}

Пайплайн построения теней для сцены с динамическими объектами при использовании карт теней сейчас устроен так:
\begin{itemize} 
    \item На входе - две карты теней. На одной отрисовываются статические объекты, на другой - динамические
    \item При изменении положения камеры идет перестройка карты теней для статических объектов
    \item Каждые 0,05 секунд или при изменении положения камеры идет перестройка карты теней для группы динамических объектов
\end{itemize}

}

\section{Section 4}
\subsection{Subsection 4.1}
\frame{{Вывод} 
{}

На данный момент:

\begin{itemize}  
\item Есть работающая песочница. Реализован классический метод каскадных теней.
\item Разработана система для разделения статических и динамических объектов.
\item Начата работа по имплементации аналитических теней.
\end{itemize}

}

\subsection*{Section 5}
\frame{{Карты теней} 

Алгоритм:
\begin{itemize} 
\item Сначала происходит построение сцены в Z-буфер с точки зрения источника света, 
чтобы установить видимые из этой точки пикселы и расстояния до них. 
\item После этого сцена перестраивается уже с точки зрения наблюдателя. 
\item Если пиксель невидим с точки зрения источника света, он закрашивается тёмным цветом.
\end{itemize}
\begin{columns}
    \column{0.3\textwidth}
    \includegraphics[width=\textwidth]{depthmap.PNG}
    \column{0.3\textwidth}
    \includegraphics[width=\textwidth]{shadows.PNG}
\end{columns}
}

\subsection*{Section 6}
\frame{{Артефакты}

\vspace{0.5em}
\begin{columns}
    \column{0.3\textwidth}
    Алиасинг
    \includegraphics[width=\textwidth]{aliasing-no-pcf.jpg}
    \column{0.3\textwidth}
    <<Теневые прыщи>>
    \includegraphics[width=\textwidth]{shadow-acne.jpg}
    \column{0.3\textwidth}
    Обрезание теней
    \includegraphics[width=\textwidth]{cutoutshadows.jpg}
\end{columns}
}

\end{document}